\section{The spin-one quantum-link Schwinger model}
\label{sec:spin1-qlm-model}

We consider the quantum-link formulation of the one-dimensional Schwinger
model with staggered fermionic matter on the sites and spin-one gauge degrees
of freedom on the links.  The fermionic annihilation and number operators at
site $n$ are denoted by $c_n$ and
$N_n=c_n^\dagger c_n$, respectively.  The electric-field operator on link
$n$ is
\begin{equation}
    E_n = S_n^z,
    \qquad E_n\in\{-1,0,+1\},
\end{equation}
and $S_n^\pm$ raise or lower the electric flux on that link.

The Hamiltonian is
\begin{equation}
\label{eq:spin1-qlm-hamiltonian}
    H =
    -t\sum_n
    \left(c_n^\dagger S_n^+c_{n+1}+\mathrm{H.c.}\right)
    +m\sum_n(-1)^nN_n
    +\frac{g^2}{2}\sum_n E_n^2 .
\end{equation}
The first term moves a fermion between adjacent sites while changing the
electric flux on the intervening link.  The second term is the staggered mass
term, and the last term is the electric-field energy.  The spin-one ladder
matrix elements are retained explicitly:
\begin{equation}
\label{eq:spin1-ladder-element}
    \langle E+1|S^+|E\rangle
    =\sqrt{S(S+1)-E(E+1)},
    \qquad S=1.
\end{equation}
Consequently, both allowed raising processes, $-1\to0$ and $0\to1$, have
matrix element $\sqrt{2}$.

Gauge invariance is generated by the local Gauss-law operators
\begin{equation}
\label{eq:gauss-law}
    G_n=E_n-E_{n-1}-N_n+q_n,
    \qquad
    q_n=\frac{1-(-1)^n}{2}.
\end{equation}
Physical states satisfy
\begin{equation}
    G_n|\Psi_{\mathrm{phys}}\rangle=0
\end{equation}
at every matter site.  Here $q_n=0$ on even sites and $q_n=1$ on odd sites.

For open boundary conditions (OBC), a chain of $L$ matter sites contains
$L-1$ dynamical links.  We denote the fixed external electric fields by
$e_{\mathrm L}=E_{-1}$ and $e_{\mathrm R}=E_{L-1}$, so that the geometry is
\begin{equation}
 e_{\mathrm L}
 \;--\;N_0\;--\;E_0\;--\;N_1\;--\;\cdots\;--\;
 E_{L-2}\;--\;N_{L-1}\;--\;e_{\mathrm R}.
\end{equation}
The hopping and electric-energy sums in Eq.~\eqref{eq:spin1-qlm-hamiltonian}
then run over the $L-1$ dynamical links.  The external fields specify a
superselection sector and are not included as dynamical degrees of freedom.
Their electric energies, if included, would give only a sector-dependent
constant.

For periodic boundary conditions (PBC), there are $L$ matter sites and $L$
dynamical links, with link $n$ joining sites $n$ and $(n+1)\bmod L$.  All site
and link indices in Eqs.~\eqref{eq:spin1-qlm-hamiltonian} and
\eqref{eq:gauss-law} are then understood modulo $L$.

Summing Gauss's law over an open chain gives
\begin{equation}
\label{eq:obc-total-number}
    N_f\equiv\sum_{n=0}^{L-1}N_n
    =e_{\mathrm R}-e_{\mathrm L}
    +\sum_{n=0}^{L-1}q_n.
\end{equation}
For even $L$, this becomes
\begin{equation}
    N_f=\frac{L}{2}+e_{\mathrm R}-e_{\mathrm L}.
\end{equation}
For PBC the electric-field differences telescope to zero, and hence
$N_f=\sum_nq_n=L/2$ for even $L$.

\section{Exact diagonalization and construction of the physical basis}
\label{sec:spin1-qlm-ed}

The exact-diagonalization calculation is performed directly in the physical
Hilbert space.  In particular, we do not construct the full tensor-product
space and subsequently impose Gauss's law through an energy penalty.  Instead,
the matter occupations and electric fields are generated so that
Eq.~\eqref{eq:gauss-law} is satisfied identically.

For OBC, Gauss's law can be rearranged into the recursion relation
\begin{equation}
\label{eq:flux-recursion}
    E_n=E_{n-1}+N_n-q_n.
\end{equation}
The construction begins with the fixed incoming field
$E_{-1}=e_{\mathrm L}$.  At site $n$, each of the two possible occupations
$N_n\in\{0,1\}$ is considered.  Equation~\eqref{eq:flux-recursion} then
determines the outgoing electric field uniquely.  At an internal site, the
branch is retained only when
\begin{equation}
    E_n\in\{-1,0,+1\},
\end{equation}
as required by the spin-one link Hilbert space.  If this condition fails, the
branch is discarded immediately.  At the final matter site, the outgoing
field is not dynamical and the branch is accepted only if it agrees with the
chosen right boundary condition,
\begin{equation}
    E_{L-2}+N_{L-1}-q_{L-1}=e_{\mathrm R}.
\end{equation}
Every branch that survives this recursive procedure is therefore a physical
basis state in the boundary-flux sector $(e_{\mathrm L},e_{\mathrm R})$.

As an example, consider $L=4$ and
$e_{\mathrm L}=e_{\mathrm R}=0$.  The background charges are
$(q_0,q_1,q_2,q_3)=(0,1,0,1)$.  The recursion produces the six states
\begin{align}
 (N_0,N_1,N_2,N_3) &=(0,0,1,1), & (E_0,E_1,E_2)&=(0,-1,0),\\
                    &=(0,1,0,1), &                    &=(0,0,0),\\
                    &=(0,1,1,0), &                    &=(0,0,1),\\
                    &=(1,0,0,1), &                    &=(1,0,0),\\
                    &=(1,0,1,0), &                    &=(1,0,1),\\
                    &=(1,1,0,0), &                    &=(1,1,1).
\end{align}
All these states have $N_f=2$, in agreement with
Eq.~\eqref{eq:obc-total-number}.

For PBC, a convenient implementation enumerates the $3^L$ electric-field
configurations and reconstructs the matter occupation from
\begin{equation}
\label{eq:pbc-matter-from-flux}
    N_n=E_n-E_{(n-1)\bmod L}+q_n.
\end{equation}
A flux configuration is retained only if the reconstructed value belongs to
$\{0,1\}$ at every site.  This produces the exact periodic physical basis.

After assigning an integer index $i$ to every physical state $|i\rangle$, the
diagonal mass and electric-field contributions are evaluated directly.  The
kinetic term connects only physical states.  In particular,
\begin{equation}
    c_n^\dagger S_n^+c_{n+1}:
    \quad (N_n,N_{n+1},E_n)
    =(0,1,E_n)\longmapsto(1,0,E_n+1),
\end{equation}
while its Hermitian conjugate implements
\begin{equation}
    (N_n,N_{n+1},E_n)
    =(1,0,E_n)\longmapsto(0,1,E_n-1).
\end{equation}
The corresponding matrix element is $-t$ multiplied by the spin ladder
factor in Eq.~\eqref{eq:spin1-ladder-element}.  For PBC, a hopping process
across the chosen fermionic ordering cut additionally carries the standard
Fock-space sign $(-1)^{N_f-1}$.

The Hamiltonian is assembled as a sparse matrix.  For small Hilbert spaces the
complete spectrum can be obtained by dense diagonalization, while for larger
systems an iterative sparse eigensolver is used to calculate the lowest
eigenvalues and eigenvectors.  Useful validation tests include checking
Hermiticity, verifying $G_n=0$ for every basis state, and confirming that each
kinetic transition maps to another state in the physical basis.
